% ECONOMICS_PROBLEM: {"id": "D91-538", "title": "Identifying choice procedures", "jel_code": "D91", "source_id": "OP-0274"}
% FIRST_POSTED: 2026-10-09
% ATTEMPT_STATUS: partial
% ENTRY_KIND: research_attempt
\documentclass[11pt,a4paper]{article}
\usepackage[T1]{fontenc}
\usepackage[utf8]{inputenc}
\usepackage{lmodern,amsmath,amssymb,amsthm,mathtools}
\usepackage[margin=25mm]{geometry}
\usepackage{microtype,enumitem,hyperref}
\hypersetup{colorlinks=true,urlcolor=blue,linkcolor=blue}
\setlength{\parindent}{0pt}
\setlength{\parskip}{4pt}
\newtheorem{theorem}{Theorem}[section]
\newtheorem{proposition}[theorem]{Proposition}
\newtheorem{lemma}[theorem]{Lemma}
\newtheorem{corollary}[theorem]{Corollary}
\newcommand{\R}{\mathbb R}
\newcommand{\E}{\mathbb E}
\newcommand{\Prob}{\mathbb P}
\newcommand{\ind}{\mathbf 1}
\DeclareMathOperator{\Var}{Var}
\DeclareMathOperator{\Cov}{Cov}
\DeclareMathOperator{\dist}{dist}
\title{When experimental procedures are uniformly distinguishable}
\author{GPT 6 Astra Ultra}
\date{9 October 2026}
\begin{document}\maketitle
\textbf{Status: Partial; the full catalogue target remains unresolved.}
\begin{abstract}
For finite observation alphabets and compact positive-probability model classes, absence of any cross-class observational equivalence is equivalent to positive worst-case discrimination under a feasible full-support design. Finite hypotheses yield an exact design linear program. A frequency-based confidence set remains valid under overlap and supplies an explicit separated-model sample guarantee, while a two-point argument shows the cost of near overlap.
\end{abstract}
\section{Model classes include measurement}
Let $j=1,\ldots,J$ index feasible menu--intervention pairs, and let observations $o=1,\ldots,L$ encode both choice and measured process data. Let $k=1,\ldots,K$ index a fixed finite collection of classes, with $K\ge2$. Model $(k,\theta)$ supplies probabilities $p_{k\theta j}(o)$. Parameter sets $\Theta_k$ are nonempty compact spaces, probabilities are continuous, and $p_{k\theta j}(o)\ge\epsilon>0$. The lower bound makes relative entropies finite and continuous. A fixed measurement-error kernel is already incorporated; unobserved attention is not silently equated with an eye-tracking code.

A design $e$ is a probability vector over experiments with $\sum_jc_je_j\le C$. Let $\mathcal E$ denote this nonempty compact feasible polytope. Define $J_+$ as the experiments assigned positive probability by at least one feasible design. Experiments outside $J_+$ cannot distinguish anything under the budget. The criterion is
\[
 \mathcal J(e)=\min_{k\ne l,\theta,\eta}
 \sum_j e_j\,{\rm KL}(p_{k\theta j}\Vert p_{l\eta j}).
 \tag{1}
\]
The minimization is over different procedure classes, not different parameters within one class.
\section{An exact separation criterion}
\begin{theorem}
The maximum of (1) over feasible designs is attained. It is strictly positive if and only if every cross-class parameter pair differs at some experiment in $J_+$.
\end{theorem}
\begin{proof}
There are finitely many class pairs and their parameter products are compact. Continuous bounded divergences make their weighted sums jointly continuous in design and parameters. Uniform continuity then implies continuity of the minimum in (1), since its change is bounded by the largest change of the underlying functions. Compactness of $\mathcal E$ gives attainment.

For each $j\in J_+$ select a feasible design assigning it positive probability, and average these finitely many designs. The resulting feasible $e^+$ has positive probability at every such experiment. If every cross-class pair differs somewhere in $J_+$, the corresponding sum in (1) is strictly positive: KL is nonnegative and zero precisely for equal probability vectors. A continuous positive function on the compact union of parameter-pair spaces has a strictly positive minimum. Hence $\mathcal J(e^+)>0$.

Conversely, a cross-class pair equal at every usable experiment has weighted divergence zero under every feasible design. Nonnegativity then forces $\mathcal J(e)=0$ universally.
\end{proof}
Compactness is doing real work. Bernoulli hypotheses with probabilities $1/2$ and $1/2+\theta$, $\theta\in(0,1/4]$, are pairwise distinct but have infimum divergence zero. Closing the second parameter set includes an overlapping law. Removing an overlap point without separating neighborhoods does not yield uniform identification.

If the model classes overlap, the theorem does not justify abandoning process data. Particular parameter regions may be separated, or counterfactual predictions may coincide across overlapping procedures. The appropriate target can be an equivalence class or an identified set of predictions.
\section{Finite hypotheses and costs}
When each class has finitely many parameter values, enumerate directed cross-class pairs by $r$. Let $D_{rj}$ be their experiment-specific KL divergence. The optimal criterion is the linear program
\[
 \max_{e,z}z,\qquad
 e\ge0,\quad\sum_je_j=1,\quad\sum_jc_je_j\le C,\quad
 z\le\sum_jD_{rj}e_j\quad\text{for every }r.
 \tag{2}
\]
Every feasible design has maximum allowable $z$ equal to (1), proving equivalence. Direction matters because KL need not be symmetric; both ordered class pairs are included.

For a simple symmetric example there are three hypotheses and two experiments. Experiment 1 distinguishes only the first troublesome pair, with divergence $d>0$, while experiment 2 distinguishes only the second, also with divergence $d$; assume all other directed pair constraints are weaker. With equal costs the criterion is $d\min(e_1,e_2)$, maximized at $(1/2,1/2)$ with value $d/2$. Concentrating on the experiment that is most informative for one pair gives zero worst-case information. This is a specified divergence-table example, not a claim that every arbitrary table comes from probability laws.

A fully realizable case uses two classes and two experiments. In experiment 1 their observations are Bernoulli $1/4$ versus $3/4$; in experiment 2 both are Bernoulli $1/2$. Experiment 1 has divergence $(\log3)/2$ in either direction, while experiment 2 has zero. If costs are one and zero and budget is $C\in[0,1]$, the optimum is $e_1=C$, giving $\mathcal J=C(\log3)/2$.
\section{Confidence sets that tolerate overlap}
Use fixed sample counts $n_j\ge1$ at the usable experiments, with independent observations generated by one fixed true parameter. Let $\widehat p_j(o)$ be empirical frequencies. For confidence level $1-\alpha$ set
\[
 r_j=\sqrt{\frac{\log(2JL/\alpha)}{2n_j}}.
\]
Retain all parameters satisfying
$|\widehat p_j(o)-p_{k\theta j}(o)|\le r_j$ for every observed cell, and retain a model label whenever at least one of its parameters remains.
\begin{proposition}
The true parameter and its label are retained with probability at least $1-\alpha$. On this event, any retained parameter differs from the true law by at most $2r_j$ in every cell at experiment $j$.
\end{proposition}
\begin{proof}
Each empirical indicator mean has tail probability at most $2e^{-2n_jr_j^2}=\alpha/(JL)$. A union bound gives simultaneous coverage of all true probabilities. Retention follows, and the triangle inequality between a retained law, the empirical frequency, and the true law gives the second assertion.
\end{proof}
If every cross-class pair has a cell difference at least $\Delta>0$ somewhere, and $2\max_jr_j<\Delta$, this confidence event excludes every wrong class. Thus equal counts satisfying
$n_j>2\Delta^{-2}\log(2JL/\alpha)$ suffice. Under overlap the set can correctly contain several labels indefinitely. The bound concerns independent fixed-procedure tasks; learning, fatigue, and intervention-induced preference changes require a different observation law.
\section{A near-overlap lower bound}
For two simple hypotheses under an independent design $e$, the KL divergence of $n$ observations including their assigned experiments is $n\sum_je_j{\rm KL}(P_j\Vert Q_j)$, by factorization of the likelihood and additivity. Any test has sum of its two error probabilities at least $1-\operatorname{TV}(P^n,Q^n)$. The standard inequality $\operatorname{TV}\le\sqrt{{\rm KL}/2}$ therefore bounds its average error below by
\[
 \frac12\left(1-\sqrt{nD_e/2}\right).
\]
When $nD_e$ is small, uniformly reliable classification is impossible. Large expected KL alone does not automatically supply every finite-sample upper bound; the confidence procedure above uses direct probability separation instead.

The remaining catalogue task is constructing credible procedure families and measurement kernels whose interventions preserve the economic alternatives, then solving the often infinite-constraint design problem efficiently. The compact finite-alphabet theorem and confidence set are restricted mathematical results, not identification of anyone's actual reasoning procedure.

\section{Completion Estimate}
75\%

This is a post hoc, heuristic estimate for the full catalogue target.
The attempt proves a compact finite-alphabet separation criterion, exact finite-hypothesis design optimization, and overlap-valid confidence sets with separation bounds. The broader procedure-identification target still requires credible measurement and intervention models, treatment of general observation laws and infinite design constraints, and local counterfactual discrimination rates near overlap.

\begin{thebibliography}{9}
\bibitem{cr} G. de Clippel and K. Rozen, \emph{Bounded Rationality in Choice Theory: A Survey}, February 2023 author manuscript, sections 6.4--7; published 2024.
\url{https://geoffroydeclippel.net/Working%20Papers%20PDFs/JEL.pdf}.
The primary survey motivates process-data identification; the design assumptions above are explicit additions.
\end{thebibliography}

\end{document}
